\documentclass[14pt,a4paper]{extarticle}

\usepackage[T2A]{fontenc}
\usepackage[utf8]{inputenc}
\usepackage[russian]{babel}
\usepackage{cmap}
\usepackage{geometry}
\usepackage{array}
\usepackage{longtable}
\usepackage{booktabs}
\usepackage{xcolor}
\usepackage{ragged2e}
\usepackage{seqsplit}
\usepackage{listings}
\usepackage{enumitem}
\usepackage{hyperref}
\usepackage{tikz}
\usetikzlibrary{arrows.meta,positioning,shapes.geometric}

\geometry{
    left=30mm,
    right=15mm,
    top=20mm,
    bottom=20mm
}

\hypersetup{
    colorlinks=true,
    linkcolor=black,
    urlcolor=blue
}

\setlength{\parindent}{1.25cm}
\setlength{\parskip}{0.2em}
\renewcommand{\arraystretch}{1.35}
\setlength{\tabcolsep}{4pt}
\sloppy

\newcolumntype{L}[1]{>{\RaggedRight\arraybackslash}p{#1}}

\lstdefinestyle{codestyle}{
    basicstyle=\ttfamily\small,
    breaklines=true,
    frame=single,
    columns=fullflexible,
    keepspaces=true,
    showstringspaces=false,
    backgroundcolor=\color{gray!5}
}

\newcommand{\code}[1]{\texttt{\detokenize{#1}}}
\newcommand{\pathcode}[1]{{\footnotesize\nolinkurl{#1}}}

\begin{document}

\begin{titlepage}
    \begin{center}
        САНКТ-ПЕТЕРБУРГСКИЙ НАЦИОНАЛЬНЫЙ\\
        ИССЛЕДОВАТЕЛЬСКИЙ УНИВЕРСИТЕТ ИТМО

        \vspace{1.5cm}

        Дисциплина: Бэк-энд разработка

        \vspace{1.5cm}

        \textbf{Отчет}

        \vspace{0.5cm}

        Домашнее задание 5

        \vspace{0.5cm}

        \textbf{Реализация межсервисного взаимодействия посредством очередей сообщений}

        \vfill

        \begin{flushright}
            Выполнил:\\
            Цатинян Артём\\
            Группа БР1.2

            \vspace{0.8cm}

            Проверил:\\
            Добряков Д. И.
        \end{flushright}

        \vfill

        Санкт-Петербург\\
        2026 г.
    \end{center}
\end{titlepage}

\section{Задача}

Тема проекта: приложение для бронирования столиков в ресторанах.

В рамках домашнего задания требовалось подключить брокер сообщений RabbitMQ или Kafka и реализовать межсервисное взаимодействие через очередь сообщений. Для текущего проекта был выбран RabbitMQ, так как он хорошо подходит для событийной интеграции между Spring Boot сервисами и не требует избыточной инфраструктуры для учебного сценария.

\section{Ход работы}

К началу работы приложение уже было разделено на микросервисы: \code{identity-service}, \code{catalog-service}, \code{booking-service} и \code{review-service}. Синхронное взаимодействие между сервисами использовалось для операций, где результат нужен сразу, например проверка пользователя или получение контекста бронирования. Для ДЗ5 был добавлен асинхронный сценарий, в котором сервис отзывов сообщает каталогу о создании нового отзыва.

\subsection{Выбор сценария}

В качестве бизнес-сценария выбрано обновление рейтинга ресторана после публикации отзыва. Это хороший пример для очереди сообщений, потому что пользовательский запрос на создание отзыва не должен напрямую зависеть от успешности пересчета агрегированного рейтинга в каталоге. Отзыв создается в \code{review-service}, после чего событие \code{ReviewCreated} публикуется в RabbitMQ. \code{catalog-service} получает это событие и обновляет таблицу агрегированной статистики рейтинга.

\begin{center}
\begin{tikzpicture}[
    node distance=2.5cm,
    service/.style={rectangle, draw, rounded corners, align=center, minimum width=3.4cm, minimum height=1cm},
    queue/.style={cylinder, draw, shape border rotate=90, aspect=0.25, align=center, minimum width=2.7cm, minimum height=1.1cm},
    db/.style={cylinder, draw, shape border rotate=90, aspect=0.25, align=center, minimum width=2.8cm, minimum height=1.1cm},
    arrow/.style={-{Latex[length=3mm]}, thick}
]
\node[service] (review) {review-service\\создание отзыва};
\node[db, below=1.4cm of review] (reviewdb) {review schema\\reviews, outbox};
\node[queue, right=3cm of review] (rabbit) {RabbitMQ\\review.created};
\node[service, right=3cm of rabbit] (catalog) {catalog-service\\listener};
\node[db, below=1.4cm of catalog] (catalogdb) {catalog schema\\rating stats};

\draw[arrow] (review) -- node[above]{publish event} (rabbit);
\draw[arrow] (rabbit) -- node[above]{consume event} (catalog);
\draw[arrow] (review) -- (reviewdb);
\draw[arrow] (catalog) -- (catalogdb);
\end{tikzpicture}
\end{center}

\subsection{Настройка RabbitMQ}

RabbitMQ подключен как инфраструктурный сервис. Для локального запуска был добавлен файл \code{docker-compose.rabbitmq.yml}, который поднимает брокер и management UI.

\begin{lstlisting}[style=codestyle]
rabbitmq:
  image: rabbitmq:3.13-management
  ports:
    - "5672:5672"
    - "15672:15672"
  environment:
    RABBITMQ_DEFAULT_USER: storage
    RABBITMQ_DEFAULT_PASS: storage
\end{lstlisting}

Для сервисов добавлены зависимости Spring AMQP. Параметры подключения вынесены в \code{application.yml}: host, port, username и password. В Docker Compose для ЛР3 эти параметры переопределяются так, чтобы сервисы обращались к брокеру по имени контейнера \code{rabbitmq}.

\subsection{Контракт события}

Общий контракт события вынесен в модуль \code{common}. Это позволяет producer и consumer использовать один и тот же DTO и одинаковые имена exchange, queue и routing key.

\begin{longtable}{|L{0.28\textwidth}|L{0.55\textwidth}|}
\hline
\textbf{Элемент} & \textbf{Значение} \\
\hline
Exchange & \code{restaurant.review.events} \\
\hline
Queue & \code{catalog.review-created} \\
\hline
Routing key & \code{review.created} \\
\hline
Событие & \code{ReviewCreatedEvent} \\
\hline
Producer & \code{review-service} \\
\hline
Consumer & \code{catalog-service} \\
\hline
\end{longtable}

Событие содержит минимальный набор данных, необходимый каталогу для обновления рейтинга:

\begin{lstlisting}[style=codestyle]
data class ReviewCreatedEvent(
    val eventId: String,
    val reviewId: Long,
    val bookingId: Long,
    val restaurantId: Long,
    val userId: Long,
    val rating: Int,
    val occurredAt: LocalDateTime
)
\end{lstlisting}

\subsection{Публикация события}

В \code{review-service} создание отзыва и запись сообщения в outbox выполняются в рамках одной бизнес-операции. Для этого в схему \code{review} добавлена таблица \code{review_outbox}. Такой подход снижает риск потери события: если отзыв успешно сохранен, событие также остается в базе и может быть опубликовано позже.

Публикация реализована отдельным компонентом \code{ReviewOutboxPublisher}. Он периодически выбирает необработанные записи outbox, преобразует payload в \code{ReviewCreatedEvent}, отправляет событие в RabbitMQ через \code{RabbitTemplate} и после успешной отправки помечает запись как обработанную.

\begin{lstlisting}[style=codestyle]
rabbitTemplate.convertAndSend(
    ReviewEvents.EXCHANGE,
    ReviewEvents.CREATED_ROUTING_KEY,
    event
)
\end{lstlisting}

\subsection{Обработка события}

В \code{catalog-service} добавлен listener \code{ReviewCreatedEventListener}. Он подписан на очередь \code{catalog.review-created} через аннотацию \code{@RabbitListener}. После получения события сервис вызывает метод репозитория \code{applyReviewCreatedEvent}.

В каталоге добавлены две таблицы:

\begin{longtable}{|L{0.34\textwidth}|L{0.49\textwidth}|}
\hline
\textbf{Таблица} & \textbf{Назначение} \\
\hline
\pathcode{restaurant_rating_stats} & Хранит агрегированный рейтинг ресторана и количество отзывов. Эти данные используются в списке ресторанов и карточке ресторана. \\
\hline
\pathcode{processed_review_events} & Хранит идентификаторы уже обработанных событий, чтобы повторная доставка из очереди не увеличивала рейтинг дважды. \\
\hline
\end{longtable}

Метод обработки сначала пытается записать \code{eventId} в таблицу \code{processed_review_events}. Если событие уже было обработано, вставка не выполняется и пересчет рейтинга пропускается. Если событие новое, сервис блокирует текущую строку статистики ресторана, пересчитывает средний рейтинг и увеличивает счетчик отзывов.

\subsection{Идемпотентность и надежность}

Очереди сообщений обычно допускают повторную доставку события. Поэтому обработчик в \code{catalog-service} сделан идемпотентным: ключом идемпотентности является \code{eventId}. Даже если RabbitMQ повторно доставит то же сообщение, таблица \code{processed_review_events} не позволит применить его второй раз.

Также используется outbox-подход на стороне \code{review-service}. Он отделяет сохранение бизнес-данных от отправки сообщения в брокер. Если RabbitMQ временно недоступен, отзыв остается в базе, outbox-запись остается необработанной, и publisher попробует отправить событие повторно при следующем запуске.

\subsection{Проверка результата}

Для проверки сценария можно поднять RabbitMQ:

\begin{lstlisting}[style=codestyle]
docker-compose -f docker-compose.rabbitmq.yml up -d
\end{lstlisting}

Затем запускаются сервисы \code{identity-service}, \code{catalog-service}, \code{booking-service} и \code{review-service}. После авторизации пользователь создает отзыв через \code{review-service}. В результате:

\begin{enumerate}[nosep]
    \item В схеме \code{review} появляется запись в \code{reviews}.
    \item В \code{review_outbox} появляется событие, которое затем помечается как обработанное.
    \item В RabbitMQ публикуется сообщение \code{ReviewCreatedEvent}.
    \item \code{catalog-service} получает событие из очереди.
    \item В схеме \code{catalog} обновляется \code{restaurant_rating_stats}.
\end{enumerate}

Проверить инфраструктуру RabbitMQ можно через web-интерфейс:

\begin{lstlisting}[style=codestyle]
http://localhost:15672
login: storage
password: storage
\end{lstlisting}

\section{Вывод}

В рамках домашнего задания было реализовано асинхронное межсервисное взаимодействие через RabbitMQ. Сервис отзывов стал producer событий о создании отзывов, а сервис каталога стал consumer этих событий и обновляет агрегированный рейтинг ресторанов. Для надежности добавлен outbox-подход, а для защиты от повторной доставки сообщений реализована идемпотентная обработка через таблицу обработанных событий.

В результате микросервисная архитектура стала менее связанной: создание отзыва не зависит напрямую от синхронного вызова каталога, а данные рейтинга обновляются через событийную интеграцию.

\end{document}
